\chapter{Methodology}
\label{ch:design}

\section{Introduction}
This chapter explains \emph{how} we approached the system: the constraints we respected, the layered architecture we chose, and the modeling artifacts (use case, structural, behavioral, deployment, and data-flow views) that make the implementation easier to reason about. The focus stays on methodology and design rationale; detailed run configuration and test evidence appear in Chapter~IV.

\section{Problem Design and Analysis}
\subsection{Design principles}
The project has the following design principles:
\begin{enumerate}[label=\arabic*.]
	\item \textbf{Modularity:} frontend is clearly separated from API, services and data layers.
	\item \textbf{Know your user:} separate actions for students and tutors handled in the UI and API.
	\item \textbf{Reliability first:} background tasks are used to ensure responsiveness.
	\item \textbf{Maintainability:} API contracts and schema validation are used.
\end{enumerate}

\subsection{Core functional actors}
The system supports three role patterns:
\begin{itemize}
\item \textbf{Student:} learns through courses, AI Tutor, routine, and community.
\item \textbf{Tutor:} manages class operations, announcements, and monitoring views.
\item \textbf{Admin-level control:} extends management and moderation operations where configured.
\end{itemize}

\section{Overall Framework and System Architecture}
The deployed architecture has four main layers:
\begin{itemize}
\item \textbf{Frontend:} React + TypeScript application for all user-facing pages.
\item \textbf{Backend API:} FastAPI service handling authentication, CRUD, and business logic.
\item \textbf{Data and Queue Layer:} PostgreSQL for persistence and Redis for cache/queue support.
\item \textbf{AI and Worker Layer:} local model integration and worker processes for long-running jobs.
\end{itemize}

\begin{figure}[htbp]
\centering
\resizebox{0.92\textwidth}{!}{%
\begin{tikzpicture}[
    node distance=10mm and 12mm,
    layer/.style={draw, rounded corners, align=center, minimum width=3.4cm, minimum height=9mm, fill=kuetgray},
    arr/.style={-{Stealth}, thick}
]
\node[layer, fill=blue!10] (frontend) {Frontend Layer\\React + TypeScript};
\node[layer, fill=green!10, below=of frontend] (api) {Backend API Layer\\FastAPI Routers + Validation};
\node[layer, fill=orange!12, below left=of api] (data) {Data Layer\\PostgreSQL + Redis};
\node[layer, fill=purple!10, below right=of api] (worker) {Worker / AI Layer\\Queue Workers + Tutor Services};

\draw[arr] (frontend) -- (api);
\draw[arr] (api) -- (data);
\draw[arr] (api) -- (worker);
\draw[arr, dashed] (worker) -- (data);
\end{tikzpicture}
}
\caption{Layered architecture of the Cognitive Copilot system.}
\label{fig:system-layer-arch}
\end{figure}

\section{Use Case Diagram}
\begin{figure}[htbp]
\centering
\resizebox{0.90\textwidth}{!}{%
\begin{tikzpicture}[
    actor/.style={draw, rounded corners, align=center, minimum width=1.9cm, minimum height=7mm, fill=blue!8},
    usecase/.style={ellipse, draw, align=center, minimum width=3.6cm, minimum height=8mm, fill=green!10},
    boundary/.style={draw, rounded corners, minimum width=10.6cm, minimum height=6.9cm},
    link/.style={-}
]
\node[boundary] (sys) at (1.7,0) {};
\node[anchor=north west, font=\small\bfseries] at (-3.4,3.25) {Cognitive Copilot};

\node[actor] (student) at (-6.6,1.8) {Student};
\node[actor] (tutor) at (-6.6,0.0) {Tutor};
\node[actor] (admin) at (-6.6,-1.8) {Admin};

\node[usecase] (auth) at (-1.2,2.2) {Login / Profile};
\node[usecase] (courses) at (-1.2,0.9) {Courses + Topics};
\node[usecase] (routine) at (-1.2,-0.4) {Routine + Attendance};
\node[usecase] (aitutor) at (-1.2,-1.7) {AI Tutor Modes};
\node[usecase] (community) at (4.6,1.6) {Community + Threads};
\node[usecase] (analytics) at (4.6,0.0) {Analytics};
\node[usecase] (notify) at (4.6,-1.6) {Notifications};

\draw[link] (student) -- (auth);
\draw[link] (student) -- (courses);
\draw[link] (student) -- (routine);
\draw[link] (student) -- (aitutor);
\draw[link] (student) -- (community);
\draw[link] (student) -- (notify);

\draw[link] (tutor) -- (courses);
\draw[link] (tutor) -- (routine);
\draw[link] (tutor) -- (community);
\draw[link] (tutor) -- (analytics);
\draw[link] (tutor) -- (notify);

\draw[link] (admin) -- (auth);
\draw[link] (admin) -- (community);
\draw[link] (admin) -- (analytics);
\end{tikzpicture}
}
\caption{Use case diagram of key actors and system interactions.}
\label{fig:use-case-diagram}
\end{figure}

\section{Core Data Domains}
The project data model is organized into practical domains:
\begin{itemize}
\item Identity and profile data
\item Course, enrollment, topic, and material data
\item Routine and attendance data
\item Community and thread interaction data
\item Notification and analytics data
\item AI tutor interaction and progress-related data
\end{itemize}

\section{Class Diagram (High Level)}
\begin{figure}[htbp]
\centering
\resizebox{0.95\textwidth}{!}{%
\begin{tikzpicture}[
    clsbox/.style={draw, rounded corners, align=left, minimum width=3.0cm, minimum height=1.1cm, fill=orange!10},
    rel/.style={-{Stealth}, thick}
]
\node[clsbox] (user) at (-4.6,1.9) {\textbf{User}\\id, name, email, role};
\node[clsbox] (course) at (-1.0,1.9) {\textbf{Course}\\id, code, name};
\node[clsbox] (topic) at (2.6,1.9) {\textbf{Topic}\\id, title, courseId};
\node[clsbox] (material) at (6.2,1.9) {\textbf{Material}\\id, topicId, type};
\node[clsbox] (enroll) at (-1.0,0.3) {\textbf{Enrollment}\\userId, courseId};
\node[clsbox] (thread) at (2.6,0.3) {\textbf{Thread}\\courseId, creatorId};
\node[clsbox] (notification) at (6.2,0.3) {\textbf{Notification}\\userId, isRead};
\node[clsbox] (progress) at (2.6,-1.3) {\textbf{TopicProgress}\\userId, topicId};

\draw[rel] (user) -- (enroll);
\draw[rel] (course) -- (enroll);
\draw[rel] (course) -- (topic);
\draw[rel] (topic) -- (material);
\draw[rel] (user) -- (thread);
\draw[rel] (course) -- (thread);
\draw[rel] (user) -- (notification);
\draw[rel] (user) -- (progress);
\draw[rel] (topic) -- (progress);
\end{tikzpicture}
}
\caption{High level class relationship diagram for core entities.}
\label{fig:class-diagram}
\end{figure}

\section{ER Diagram (Conceptual)}
\begin{figure}[htbp]
\centering
\resizebox{0.86\textwidth}{!}{%
\begin{tikzpicture}[
    ent/.style={draw, rectangle, align=center, minimum width=2.7cm, minimum height=8mm, fill=cyan!10},
    rel/.style={-{Stealth}, thick}
]
\node[ent] (u)  at (-5.2, 2.2) {User};
\node[ent] (c)  at (-1.5, 2.2) {Course};
\node[ent] (e)  at ( 2.2, 2.2) {Enrollment};
\node[ent] (t)  at (-1.5, 0.6) {Topic};
\node[ent] (m)  at ( 2.2, 0.6) {Material};
\node[ent] (tp) at (-5.2, 0.6) {TopicProgress};
\node[ent] (th) at (-1.5,-1.0) {Thread};
\node[ent] (nt) at ( 2.2,-1.0) {Notification};

\draw[rel] (u) -- node[above, font=\scriptsize]{enrolls} (e);
\draw[rel] (c) -- node[above, font=\scriptsize]{has} (e);
\draw[rel] (c) -- node[left,  font=\scriptsize]{has} (t);
\draw[rel] (t) -- node[above, font=\scriptsize]{contains} (m);
\draw[rel] (u) -- node[left,  font=\scriptsize]{tracks} (tp);
\draw[rel] (t) -- node[above, font=\scriptsize]{for} (tp);
\draw[rel] (u) -- node[below, font=\scriptsize]{creates} (th);
\draw[rel] (c) -- node[right, font=\scriptsize]{context} (th);
\draw[rel] (u) -- node[below, font=\scriptsize]{receives} (nt);
\end{tikzpicture}
}
\caption{Conceptual ER diagram of the main project core tables.}
\label{fig:er-diagram}
\end{figure}
\section{Module Interaction Flow}
At run time the project has a uniform flow:
\begin{enumerate}[label=\arabic*.]
	\item API requests are made from the frontend pages.
	\item API validates data against schemas.
	\item Services enforce business rules and authentication.
	\item Data layer saves and returns normalized data.
	\item Frontend is kept in-sync with query cache and real-time updates.
\end{enumerate}

\begin{figure}[htbp]
\centering
\resizebox{0.96\textwidth}{!}{%
\begin{tikzpicture}[
    node distance=8mm and 10mm,
    flowstep/.style={draw, rounded corners, align=center, minimum width=3.0cm, minimum height=8mm, fill=blue!8},
    arr/.style={-{Stealth}, thick}
]
\node[flowstep] (s1) {User Interaction\\(Page action)};
\node[flowstep, right=of s1] (s2) {API Request\\(HTTP call)};
\node[flowstep, right=of s2] (s3) {Validation + Auth\\(Backend checks)};
\node[flowstep, right=of s3] (s4) {Service Logic\\(CRUD rules)};
\node[flowstep, right=of s4] (s5) {Database Update\\or Query Result};

\draw[arr] (s1) -- (s2);
\draw[arr] (s2) -- (s3);
\draw[arr] (s3) -- (s4);
\draw[arr] (s4) -- (s5);
\draw[arr, dashed] (s5.south) to[out=-90, in=-90] node[below, font=\scriptsize] {normalized response + cache refresh} (s1.south);
\end{tikzpicture}
}
\caption{End to end interaction flow from user action to synchronized UI state.}
\label{fig:module-flow}
\end{figure}

\section{Activity Diagram (Authentication Flow)}
\begin{figure}[htbp]
\centering
\resizebox{0.82\textwidth}{!}{%
\begin{tikzpicture}[
    node distance=8mm,
    activity/.style={draw, rounded corners, align=center, minimum width=3.5cm, minimum height=8mm, fill=purple!10},
    decision/.style={diamond, draw, align=center, minimum width=2.1cm, minimum height=8mm, fill=yellow!15},
    arr/.style={-{Stealth}, thick}
]
\node[activity] (a1) {User submits credentials};
\node[activity, below=of a1] (a2) {Request validation};
\node[decision, below=of a2, yshift=-1mm] (d1) {Valid?};
\node[activity, below left=of d1, xshift=-8mm] (a3) {Return error};
\node[activity, below right=of d1, xshift=8mm] (a4) {Verify user + password};
\node[decision, below=of a4, yshift=-1mm] (d2) {Authorized?};
\node[activity, below left=of d2, xshift=-8mm] (a5) {Return unauthorized};
\node[activity, below right=of d2, xshift=8mm] (a6) {Issue token + response};

\draw[arr] (a1) -- (a2);
\draw[arr] (a2) -- (d1);
\draw[arr] (d1) -- node[above, font=\scriptsize]{No} (a3);
\draw[arr] (d1) -- node[above, font=\scriptsize]{Yes} (a4);
\draw[arr] (a4) -- (d2);
\draw[arr] (d2) -- node[above, font=\scriptsize]{No} (a5);
\draw[arr] (d2) -- node[above, font=\scriptsize]{Yes} (a6);
\end{tikzpicture}
}
\caption{Activity diagram for the authentication request lifecycle.}
\label{fig:activity-auth}
\end{figure}

\section{AI Tutor Methodology}
The AI Tutor is integrated as a module of the LMS rather than a separate chatbot. It supports multiple learning modes
through a shared workspace and course/topic context. Retrieval based answering and quiz workflows are connected with
existing course entities so interactions remain academically relevant.

\section{Sequence Diagram (AI Tutor Request)}
\begin{figure}[htbp]
\centering
\resizebox{0.92\textwidth}{!}{%
\begin{tikzpicture}[x=1.05cm,y=0.9cm, font=\small]
\draw[thick] (0,0) -- (0,-5.8) node[below] {User};
\draw[thick] (3,0) -- (3,-5.8) node[below] {Frontend};
\draw[thick] (6,0) -- (6,-5.8) node[below] {Backend};
\draw[thick] (9,0) -- (9,-5.8) node[below] {AI Service};
\draw[thick] (12,0) -- (12,-5.8) node[below] {Database};

\draw[-{Stealth}] (0,-0.7) -- (3,-0.7) node[midway,above,font=\scriptsize] {Ask question};
\draw[-{Stealth}] (3,-1.3) -- (6,-1.3) node[midway,above,font=\scriptsize] {POST /ai-tutor/...};
\draw[-{Stealth}] (6,-1.9) -- (12,-1.9) node[midway,above,font=\scriptsize] {Fetch context};
\draw[-{Stealth}] (12,-2.5) -- (6,-2.5) node[midway,above,font=\scriptsize] {Context result};
\draw[-{Stealth}] (6,-3.1) -- (9,-3.1) node[midway,above,font=\scriptsize] {Generate response};
\draw[-{Stealth}, dashed] (9,-3.7) -- (6,-3.7) node[midway,above,font=\scriptsize] {Stream tokens};
\draw[-{Stealth}, dashed] (6,-4.3) -- (3,-4.3) node[midway,above,font=\scriptsize] {API stream};
\draw[-{Stealth}, dashed] (3,-4.9) -- (0,-4.9) node[midway,above,font=\scriptsize] {Render reply};
\end{tikzpicture}
}
\caption{Sequence diagram for AI Tutor request and response flow.}
\label{fig:sequence-ai-tutor}
\end{figure}

\section{Asynchronous Processing Methodology}
For tasks that can be expensive (e.g., material ingestion and OCR), asynchronous worker execution is used. This prevents
HTTP request blocking, improves responsiveness, and allows safer retry and failure handling patterns.

\begin{figure}[htbp]
\centering
\resizebox{0.88\textwidth}{!}{%
\begin{tikzpicture}[
    node distance=9mm and 12mm,
    n/.style={draw, rounded corners, align=center, minimum width=3.1cm, minimum height=8mm, fill=green!8},
    q/.style={draw, rounded corners, align=center, minimum width=3.1cm, minimum height=8mm, fill=orange!10},
    arr/.style={-{Stealth}, thick}
]
\node[n] (req) {User Upload / Action};
\node[q, right=of req] (enqueue) {Enqueue Background Job};
\node[n, right=of enqueue] (worker) {Worker Processing};
\node[n, below=of worker] (store) {Persist Result / Status};
\node[n, left=of store] (notify) {Notify UI Refresh};

\draw[arr] (req) -- (enqueue);
\draw[arr] (enqueue) -- (worker);
\draw[arr] (worker) -- (store);
\draw[arr] (store) -- (notify);
\draw[arr, dashed] (notify) -- (req);
\end{tikzpicture}
}
\caption{Asynchronous job flow used for heavy operations.}
\label{fig:async-flow}
\end{figure}

\section{Component Diagram}
\begin{figure}[htbp]
\centering
\resizebox{0.78\textwidth}{!}{%
\begin{tikzpicture}[
    node distance=9mm,
    comp/.style={draw, rounded corners, align=center, minimum width=4.6cm, minimum height=8mm, fill=gray!15},
    arr/.style={-{Stealth}, thick}
]
\node[comp, fill=blue!10] (web) {Web Client Components};
\node[comp, fill=green!10, below=of web] (api) {API Router Components};
\node[comp, fill=orange!12, below=of api] (service) {Domain Service Components};
\node[comp, fill=purple!10, below left=8mm and 7mm of service] (worker) {Worker Components};
\node[comp, fill=cyan!10, below right=8mm and 7mm of service] (data) {Database Components};
\node[comp, fill=yellow!15, left=of api, xshift=-16mm] (socket) {Realtime Socket Components};

\draw[arr] (web) -- (api);
\draw[arr] (api) -- (service);
\draw[arr] (service) -- (worker);
\draw[arr] (service) -- (data);
\draw[arr] (socket) -- (api);
\draw[arr, dashed] (worker) -- (data);
\draw[arr, dashed] (data) to[bend left=15] (api);
\end{tikzpicture}
}
\caption{Component level view of the major software blocks with interaction paths.}
\label{fig:component-diagram}
\end{figure}

\section{Deployment Diagram}
\begin{figure}[htbp]
\centering
\resizebox{0.92\textwidth}{!}{%
\begin{tikzpicture}[
    nodebox/.style={draw, rounded corners, align=center, minimum width=3.0cm, minimum height=9mm, fill=green!8},
    arr/.style={-{Stealth}, thick}
]
\node[nodebox] (browser) at (0,2.5) {User Browser};
\node[nodebox] (frontend) at (3.8,2.5) {Frontend App};
\node[nodebox] (backend) at (7.6,2.5) {Backend API};
\node[nodebox] (postgres) at (7.6,1.0) {PostgreSQL};
\node[nodebox] (redis) at (11.2,1.0) {Redis};
\node[nodebox] (ai) at (11.2,2.5) {AI Runtime / Workers};

\draw[arr] (browser) -- (frontend);
\draw[arr] (frontend) -- (backend);
\draw[arr] (backend) -- (postgres);
\draw[arr] (backend) -- (redis);
\draw[arr] (backend) -- (ai);
\draw[arr, dashed] (ai) -- (redis);
\end{tikzpicture}
}
\caption{Deployment diagram for local/full-stack runtime setup.}
\label{fig:deployment-diagram}
\end{figure}

\section{Data Flow Diagram}
\begin{figure}[htbp]
\centering
\resizebox{0.92\textwidth}{!}{%
\begin{tikzpicture}[
    ext/.style={draw, rounded corners, align=center, minimum width=2.3cm, minimum height=7mm, fill=blue!10},
    proc/.style={draw, ellipse, align=center, minimum width=2.6cm, minimum height=9mm, fill=orange!12},
    store/.style={draw, rectangle, align=center, minimum width=2.8cm, minimum height=7mm, fill=green!10},
    arr/.style={-{Stealth}, thick}
]
\node[ext] (student) at (0,1.8) {Student / Tutor};
\node[proc] (frontend) at (3.4,1.8) {Frontend Process};
\node[proc] (api) at (6.8,1.8) {Backend Process};
\node[store] (db) at (10.2,2.5) {PostgreSQL Store};
\node[store] (cache) at (10.2,1.1) {Redis / Queue Store};
\node[proc] (worker) at (6.8,0.4) {Worker Process};

\draw[arr] (student) -- node[above, font=\scriptsize]{inputs} (frontend);
\draw[arr] (frontend) -- node[above, font=\scriptsize]{requests} (api);
\draw[arr] (api) -- node[above, font=\scriptsize]{CRUD} (db);
\draw[arr] (api) -- node[below, font=\scriptsize]{enqueue} (cache);
\draw[arr] (cache) -- (worker);
\draw[arr] (worker) -- (db);
\draw[arr, dashed] (api) to[bend left=20] node[below, font=\scriptsize]{responses} (frontend);
\end{tikzpicture}
}
\caption{Data flow diagram showing request, processing, and persistence paths.}
\label{fig:data-flow-diagram}
\end{figure}

\section{Conclusion}
The general approach is implementation focused and system focused: design and implement each module with contracts, keep role
and validation rules clear, and connect all modules with stable interfaces so that key academic processes work
consistently end to end.